% Generated from locale/tr/content/computability/computability-theory/non-comp-set.tex; do not hand-edit.


\olfileid[tr]{cmp}{thy}{ncp}
\olsection{Hesaplanabilir Olmayan Kümeler Vardır}

Yukarıda her hesaplanabilir kümenin hesaplanabilir biçimde
numaralandırılabilir olduğunu gördük. Tersi doğru mudur? Aşağıdaki sonuç
genel olarak doğru olmadığını gösterir.

\begin{thm}\ollabel{thm:K-0}
$K_0$, $\Setabs{\tuple{e,x}}{\cfind{e}(x)\fdefined}$ kümesi olsun.
O zaman $K_0$ hesaplanabilir biçimde numaralandırılabilir, fakat
hesaplanabilir değildir.
\end{thm}

\begin{proof}
$K_0$'ın hesaplanabilir biçimde numaralandırılabilir olduğunu görmek için,
onun
\[
  f(z)=\umin{y}{(\len{z}=2\land T((z)_0,(z)_1,y))}
\]
ile tanımlanan $f$ fonksiyonunun tanım kümesi olduğuna dikkat edin.
Gerçekten $\cfind{e}(x)$ tanımlıysa $f(\tuple{e,x})$ duran bir hesap dizisi
bulur; $\cfind{e}(x)$ tanımsızsa $f(\tuple{e,x})$ de tanımsızdır. $z$ bir
çifti bile kodlamıyorsa $f(z)$ yine tanımsızdır.

$K_0$'ın hesaplanabilir olmaması tam olarak durma probleminin karar
verilemezliğidir; bkz. \olref[hlt]{thm:halting-problem}.
\end{proof}

$K_0$, $\cfind{e}(x)\fdefined$ olacak $\tuple{e,x}$ çiftleri kümesidir;
yani $\tuple{e,x}\in K_0$ ancak ve ancak $\cfind{e}$, $x$ girdisinde
tanımlıysa, başka bir deyişle duruyorsa geçerlidir. Bu nedenle $K_0$'a
``durma kümesi'' de denir. $K=\Setabs{e}{\cfind{e}(e)\fdefined}$ ise
``kendi üzerinde durma kümesi''dir ve çoğu kez örnek karar verilemez küme
olarak kullanılır.

\begin{thm}\ollabel{thm:K}
Kendi üzerinde durma kümesi
$K=\Setabs{e}{\cfind{e}(e)\fdefined}$ !!{c.e.} fakat karar verilebilir
değildir.
\end{thm}

\begin{proof}
  $K$ karar verilebilir, yani karakteristik fonksiyonu $\Char{K}$
  hesaplanabilir olsun. Şöyle tanımlayalım:
  \[
    d(e)=\begin{cases}
      1 & \text{$\Char{K}(e)=0$ ise}\\
      \fundefined & \text{aksi durumda.}
    \end{cases}
  \]
  $k$, $d$'nin indisi, yani $d\simeq\cfind{k}$ olsun. O zaman
  $d(k)\simeq\cfind{k}(k)$ olur. Bu, $d$ tanımından çıkan
  $d(k)\fdefined$ ancak ve ancak $\cfind{k}(k)\fundefined$ olgusuyla
  çelişir.

  $K$, $f(x)=\umin{y}{T(x,x,y)}$ fonksiyonunun tanım kümesidir;
  dolayısıyla !!{c.e.}
\end{proof}

